\begin{abstract}
\noindent This paper examines whether satellite-predicted crop yields causally affect violent conflict at the subnational level across seven Sub-Saharan African countries---Burkina Faso, Ethiopia, Mali, Malawi, Niger, Nigeria, and Tanzania---over the period 2010--2024. A central identification challenge is that observed yields are endogenous to conflict: armed violence disrupts planting, cultivation, and harvest. To address this, I use an XGBoost model trained on GROW-Africa administrative maize yield records and 210 satellite and environmental features to generate predicted yields for 1,338 GADM Admin-2 districts annually. These predictions serve as the key regressor in a two-way fixed-effects panel with Admin-2 and year fixed effects. The preferred TWFE OLS log-log specification yields a negative and highly significant elasticity: a 10\% increase in ML-predicted maize yield is associated with a 6.7\% reduction in log conflict events in the 3-month post-harvest window ($\hat{\beta} = -0.666$, $p < 0.001$), robust to spatial lag control and longer time horizons. The paper's most novel contribution is a systematic reversal of this relationship across agro-ecological zones: in Non-Sahel countries (Ethiopia, Malawi, Nigeria, Tanzania) the yield effect is negative ($\hat{\beta} = -0.467$, $p < 0.001$), consistent with an opportunity-cost mechanism, while in Sahel countries (Burkina Faso, Mali, Niger) higher predicted yields are associated with \textit{more} conflict (net elasticity $= +2.383$), consistent with a resource-predation channel in which armed groups appropriate agricultural surplus. The interaction of yield with Sahel zone membership is $+3.513$ ($p < 0.001$). These findings demonstrate that aggregate estimates mask a fundamental heterogeneity in the direction of the yield--conflict relationship and carry direct implications for the design of agricultural interventions in conflict-affected settings.
\end{abstract}
